\subsection{Minimum Edit Distance}
\label{sec:minimum-edit-distance}
\slidesrc{Syntax: RE \& word: 66--82}

\subsubsection{How similar are two strings?}
\slidesrc{Syntax: RE \& word: 66--68, 72--73}

Many tasks need a measure of how close two strings are:
\begin{itemize}
  \item \textbf{Spell correction}: the user typed \emph{graffe}. Which is closest:
        \emph{graf}, \emph{graft}, \emph{grail} or \emph{giraffe}?
  \item \textbf{Computational biology}: align two sequences of nucleotides (the bases
        Adenine, Guanine, Thymine/Uracil and Cytosine of DNA/RNA), giving each letter a
        partner letter or a gap \texttt{-}:
\begin{center}
\small\ttfamily
\begin{tabular}{@{}l@{}}
AGGCTATCACCTGACCTCCAGGCCGATGCCC\\
TAGCTATCACGACCGCGGTCGATTTGCCCGAC\\[3pt]
\textrm{\normalsize aligned:}\\
-AGGCTATCACCTGACCTCCAGGCCGA--TGCCC---\\
TAG-CTATCAC--GACCGC--GGTCGATTTGCCCGAC
\end{tabular}
\end{center}
        Alignments are used in molecular phylogenetics (evolutionary relationships from
        hereditary differences in DNA), in comparing the genomes of ill and healthy
        individuals, and to understand the function of unknown proteins.
  \item \textbf{Evaluating machine translation and speech recognition}: compare the
        system's output (hypothesis H) with a reference (R), counting substitutions (S),
        insertions (I) and deletions (D):
\begin{center}
\small
\begin{tabular}{@{}l *{9}{c}@{}}
R: & Spokesman & confirms & *   & senior & government & adviser & was & shot & *    \\
H: & Spokesman & said     & the & senior & *          & adviser & was & shot & dead \\
   &           & S        & I   &        & D          &         &     &      & I
\end{tabular}
\end{center}
  \item \textbf{Named entity extraction and coreference}: deciding that \emph{IBM Inc.}
        and \emph{IBM}, or \emph{Stanford President John Hennessy} and \emph{Stanford
        University President John Hennessy}, refer to the same entity.
\end{itemize}

\subsubsection{Definition}
\slidesrc{Syntax: RE \& word: 69--71}

\begin{nbdef}{Minimum edit distance}
The \textbf{minimum edit distance} between two strings is the minimum number of
editing operations
\begin{itemize}[nosep]
  \item \textbf{insertion} of a character,
  \item \textbf{deletion} of a character,
  \item \textbf{substitution} of a character with another,
\end{itemize}
needed to transform one string into the other. If operations have \emph{costs}, it is the
minimum total cost. With cost 1 for insertions and deletions and \textbf{cost 2 for
substitutions} it is called the \textbf{Levenshtein distance} (the version used in these
notes unless stated otherwise).
\end{nbdef}

An \textbf{alignment} shows the operations explicitly. For \emph{intention} and
\emph{execution} (\texttt{*} = gap):
\[
\begin{array}{*{10}{c}}
  \texttt{I} & \texttt{N} & \texttt{T} & \texttt{E} & \texttt{*} & \texttt{N} & \texttt{T} & \texttt{I} & \texttt{O} & \texttt{N}\\
  | & | & | & | & | & | & | & | & | & |\\
  \texttt{*} & \texttt{E} & \texttt{X} & \texttt{E} & \texttt{C} & \texttt{U} & \texttt{T} & \texttt{I} & \texttt{O} & \texttt{N}\\
  \text{d} & \text{s} & \text{s} & & \text{i} & \text{s} & & & &
\end{array}
\]
One deletion (d), three substitutions (s) and one insertion (i):
\begin{itemize}[nosep]
  \item if every operation costs 1, the distance is $1 + 3 + 1 = 5$;
  \item with Levenshtein costs (substitution $= 2$), it is $1 + 3\cdot 2 + 1 = 8$.
\end{itemize}

\begin{nbexample}{the spelling question}
Levenshtein distances from \emph{graffe}: \emph{giraffe} $=1$ (insert \texttt{i}),
\emph{graf} $=2$ (delete \texttt{f} and \texttt{e}), \emph{graft} $=3$ (delete
\texttt{f}, substitute \texttt{e}$\to$\texttt{t}), \emph{grail} $=5$. So \emph{giraffe}
is the best correction. (With all costs equal to 1: $1, 2, 2, 3$.)
\end{nbexample}

\subsubsection{Edit distance as search}
\slidesrc{Syntax: RE \& word: 74--78}

Finding the minimum edit distance is a search for a path (a sequence of edits) from the
start string to the final string:
\begin{itemize}[nosep]
  \item \textbf{initial state}: the word we are transforming;
  \item \textbf{operators}: insert, delete, substitute;
  \item \textbf{goal state}: the word we are trying to reach;
  \item \textbf{path cost}: what we minimise, the number (or cost) of edits.
\end{itemize}
The space of all edit sequences is huge, so we cannot explore it naively. But many
distinct paths end up in the \emph{same} state (the same prefix of $X$ turned into the
same prefix of $Y$), and for each such state we only need to remember the
\emph{cheapest} way to reach it. That is exactly the idea of dynamic programming.

\begin{nbdef}{The subproblems}
For strings $X$ of length $n$ and $Y$ of length $m$, let
\[
  D(i, j) = \text{edit distance between } X[1..i] \text{ and } Y[1..j],
\]
i.e.\ between the first $i$ characters of $X$ and the first $j$ characters of $Y$. The
edit distance between $X$ and $Y$ is $D(n, m)$.
\end{nbdef}

\textbf{Dynamic programming} computes $D(n,m)$ \emph{bottom-up} in a table, solving the
problem by combining solutions to subproblems: first $D(i,j)$ for small $i, j$, then
larger ones from the smaller values already computed.

\subsubsection{The algorithm}
\slidesrc{Syntax: RE \& word: 79--80}

\begin{nbfact}{Levenshtein distance by dynamic programming}
\textbf{Initialisation} (distance from the empty string):
\[ D(i, 0) = i, \qquad D(0, j) = j. \]
\textbf{Recurrence}: for each $i = 1, \dots, n$ and each $j = 1, \dots, m$,
\[
  D(i,j) = \min \begin{cases}
     D(i-1, j) + 1 & \text{deletion of } X[i]\\
     D(i, j-1) + 1 & \text{insertion of } Y[j]\\
     D(i-1, j-1) + \begin{cases} 2 & \text{if } X[i] \neq Y[j]\\ 0 & \text{if } X[i] = Y[j]\end{cases}
       & \text{substitution or copy}
  \end{cases}
\]
\textbf{Termination}: $D(n, m)$ is the distance.
\end{nbfact}

\begin{nbintuition}[reading the recurrence]
To turn $X[1..i]$ into $Y[1..j]$, look at the \emph{last} column of an optimal
alignment. Either the last character of $X$ was deleted (and the rest is an optimal
alignment of $X[1..i-1]$ with $Y[1..j]$), or the last character of $Y$ was inserted, or
the two last characters are aligned with each other (free if equal, a substitution
otherwise). The minimum over the three cases is the answer.
\end{nbintuition}

The general version (Jurafsky and Martin) uses arbitrary costs \texttt{del-cost},
\texttt{ins-cost}, \texttt{sub-cost}, which can depend on the letters, to model that some
letters are more likely to be inserted or confused than others; copying a letter costs
$\texttt{sub-cost}(x,x) = 0$.

\begin{nbwarning}[indices on the slide]
The slide defines $X$ of length $n$ and $Y$ of length $m$, but the recurrence loops over
$i = 1\dots M$, $j = 1\dots N$ and terminates with $D(N, M)$. The letters are swapped:
$i$ indexes $X$ and runs to $n$, $j$ indexes $Y$ and runs to $m$, and the result is
$D(n,m)$. The bottom-up loop also covers $1 \le i \le n$, $1 \le j \le m$ (the slide
writes strict inequalities).
\end{nbwarning}

\subsubsection{The table and the backtrace}
\slidesrc{Syntax: RE \& word: 81--82}

\begin{nbexample}{intention $\to$ execution}
Rows: the source \emph{intention}; columns: the target \emph{execution}; \texttt{\#} is
the empty prefix. Each cell is computed from its left, upper and upper-left neighbours.
\begin{center}
\small
\begin{tabular}{@{}c|cccccccccc@{}}
 & \# & \textbf{e} & \textbf{x} & \textbf{e} & \textbf{c} & \textbf{u} & \textbf{t} & \textbf{i} & \textbf{o} & \textbf{n} \\ \hline
\# & \cellcolor{nbteal!25}\textbf{0} & 1 & 2 & 3 & 4 & 5 & 6 & 7 & 8 & 9 \\
\textbf{i} & \cellcolor{nbteal!25}\textbf{1} & 2 & 3 & 4 & 5 & 6 & 7 & 6 & 7 & 8 \\
\textbf{n} & 2 & \cellcolor{nbteal!25}\textbf{3} & 4 & 5 & 6 & 7 & 8 & 7 & 8 & 7 \\
\textbf{t} & 3 & 4 & \cellcolor{nbteal!25}\textbf{5} & 6 & 7 & 8 & 7 & 8 & 9 & 8 \\
\textbf{e} & 4 & 3 & 4 & \cellcolor{nbteal!25}\textbf{5} & \cellcolor{nbteal!25}\textbf{6} & 7 & 8 & 9 & 10 & 9 \\
\textbf{n} & 5 & 4 & 5 & 6 & 7 & \cellcolor{nbteal!25}\textbf{8} & 9 & 10 & 11 & 10 \\
\textbf{t} & 6 & 5 & 6 & 7 & 8 & 9 & \cellcolor{nbteal!25}\textbf{8} & 9 & 10 & 11 \\
\textbf{i} & 7 & 6 & 7 & 8 & 9 & 10 & 9 & \cellcolor{nbteal!25}\textbf{8} & 9 & 10 \\
\textbf{o} & 8 & 7 & 8 & 9 & 10 & 11 & 10 & 9 & \cellcolor{nbteal!25}\textbf{8} & 9 \\
\textbf{n} & 9 & 8 & 9 & 10 & 11 & 12 & 11 & 10 & 9 & \cellcolor{nbteal!25}\textbf{8} \\
\end{tabular}
\captionof{table}{Levenshtein table. The bottom-right cell is the distance (8). Shaded
cells: one minimum-cost path found by the backtrace.}
\end{center}
Some cells, step by step:
\begin{align*}
  D(1,1) &= \min\bigl(D(0,1)+1,\ D(1,0)+1,\ D(0,0)+2\bigr) = \min(2, 2, 2) = 2\\
  D(4,1) &= \min\bigl(D(3,1)+1,\ D(4,0)+1,\ D(3,0)+0\bigr) = \min(5, 5, 3) = 3\\
  D(9,9) &= \min\bigl(D(8,9)+1,\ D(9,8)+1,\ D(8,8)+0\bigr) = \min(10, 10, 8) = 8
\end{align*}
In $D(1,1)$ the letters differ (\texttt{i}, \texttt{e}), so the diagonal costs 2; in
$D(4,1)$ and $D(9,9)$ they are equal, so the diagonal is a free copy.
\end{nbexample}

\textbf{Backtrace.} The table gives the distance, not the alignment. To get the
alignment, when filling each cell we remember \emph{which} neighbour(s) gave the minimum
(an arrow $\leftarrow$, $\uparrow$ or $\nwarrow$). After the table is full, we start from
the bottom-right cell and follow the arrows back to $(0,0)$: $\uparrow$ is a deletion,
$\leftarrow$ an insertion, $\nwarrow$ a substitution or a copy. The shaded path, read from
the top-left, is: delete \texttt{i}; \texttt{n}$\to$\texttt{e}; \texttt{t}$\to$\texttt{x};
copy \texttt{e}; insert \texttt{c}; \texttt{n}$\to$\texttt{u}; copy \texttt{t},
\texttt{i}, \texttt{o}, \texttt{n}. That is exactly the alignment shown above, with cost
$1 + 2 + 2 + 0 + 1 + 2 = 8$. When several neighbours tie, there are several optimal
alignments.

\begin{nbintuition}[cost of the algorithm]
The table has $(n+1)(m+1)$ cells, each computed in constant time from three neighbours,
so the algorithm takes $O(nm)$ time and space; the backtrace adds $O(n+m)$. The naive
search over all edit sequences is exponential.
\end{nbintuition}

\subsubsection{Exercises}
\notinslides

\begin{nbexercise}{1}{small tables}
Compute the Levenshtein table for \emph{cat} $\to$ \emph{cut} and for \emph{cat} $\to$
\emph{act}. Give the distance and an optimal alignment for each.
\end{nbexercise}
\begin{nbsolution}
\emph{cat} $\to$ \emph{cut} (rows c, a, t; columns c, u, t):
\[
\begin{array}{c|cccc}
   & \# & c & u & t\\ \hline
\# & 0 & 1 & 2 & 3\\
c  & 1 & 0 & 1 & 2\\
a  & 2 & 1 & 2 & 3\\
t  & 3 & 2 & 3 & 2
\end{array}
\]
Distance 2: copy \texttt{c}, substitute \texttt{a}$\to$\texttt{u} (cost 2), copy
\texttt{t}. Deleting \texttt{a} and inserting \texttt{u} also costs 2: there are two
optimal alignments.

\emph{cat} $\to$ \emph{act} (columns a, c, t):
\[
\begin{array}{c|cccc}
   & \# & a & c & t\\ \hline
\# & 0 & 1 & 2 & 3\\
c  & 1 & 2 & 1 & 2\\
a  & 2 & 1 & 2 & 3\\
t  & 3 & 2 & 3 & 2
\end{array}
\]
Distance 2, e.g.\ delete \texttt{c}, copy \texttt{a}, insert \texttt{c}, copy \texttt{t}.

\emph{Lesson}: swapping two adjacent letters costs 2. The Damerau--Levenshtein variant
adds \emph{transposition} as a fourth operation, because swaps are a very common typing
error.
\end{nbsolution}

\begin{nbsummary}[syntax at the word level]
\begin{itemize}[nosep]
  \item Word segmentation: spaces are only an approximation; Chinese, Japanese, Thai have
        no word spaces. MaxMatch = greedy longest dictionary word (good for Chinese,
        bad for English).
  \item Stop words: very frequent, little content; no universal list. A word frequent in
        a document may be topical: compare with the corpus.
  \item $\text{tf}(t,d)=\#t/\text{len}(d)$,
        $\text{idf}(t,C)=\log\bigl(\abs{C}/\abs{\{d: t\in d\}}\bigr)$,
        tf-idf = tf $\cdot$ idf; a word in every document has idf $=0$.
  \item Stemming = heuristic suffix chopping (Porter: measure $m$ in $[C](VC)^m[V]$,
        conditional rules in 7 steps; over- and under-stemming). Lemmatisation =
        dictionary + morphology, returns the lemma.
  \item Minimum edit distance: insertions, deletions, substitutions; Levenshtein:
        substitution costs 2. Dynamic programming: $D(i,0)=i$, $D(0,j)=j$,
        $D(i,j)$ = minimum over the upper neighbour $+1$ (deletion), the left
        neighbour $+1$ (insertion), the diagonal neighbour $+2$ (substitution) or $+0$
        (copy); the backtrace gives the alignment; time and space $O(nm)$.
\end{itemize}
\end{nbsummary}
